\documentclass[11pt,a4paper]{article}
\usepackage[margin=1in]{geometry}
\usepackage{times}
\usepackage{amsmath}

\linespread{1.15}
\setlength{\parskip}{6pt}
\setlength{\parindent}{0pt}

\title{\textbf{Spectrum Sensing and Cognitive IoT Network}\\[4pt]
\large Cooperative Energy-Detection-Based Spectrum Sensing: A Python Simulation Study}
\author{Siva Dheeraj Ethamsetti \\ Phani Sai Nadha Reddy}
\date{}

\begin{document}
\maketitle

\begin{abstract}
The radio spectrum used by IoT deployments is increasingly congested, while licensed
bands can remain idle for long periods. Cognitive radio allows IoT devices to
use a licensed band opportunistically, provided they can reliably detect whether the
primary (licensed) user is active. This project designs and evaluates an energy
detector for secondary IoT nodes and studies how cooperative sensing improves
detection reliability. Each node measures the received signal energy over $N = 1000$
samples of a BPSK primary signal in additive white Gaussian noise and compares it
with a threshold fixed for a false-alarm probability of $0.1$; several nodes send
1-bit decisions to a fusion center that applies the OR rule. A Monte Carlo
simulation in Python (2000 trials per point, SNR swept from $-20$ to $+5$~dB) is
validated against the Gaussian-approximation theory and used to produce
probability-of-detection, ROC, and cooperative-gain curves. At $-10$~dB, the
detection probability rises from about $0.80$ for a single node to $0.99$ with three
nodes and $1.0$ with five nodes, at the cost of a higher overall false-alarm rate
under the OR rule. The results quantify the reliability gain that cooperation brings
to low-SNR spectrum sensing in cognitive IoT networks.
\end{abstract}

\section{Introduction}

The number of connected IoT devices keeps growing, and the unlicensed bands most of
them share are increasingly crowded. At the same time, licensed bands can remain
idle for long stretches: when they do, the scarcity is not of spectrum itself, but
of permission to use it. Cognitive radio addresses this
mismatch by letting unlicensed (secondary) users operate in a licensed band whenever
the primary user is absent, and vacating it when the primary user returns.

Opportunistic access is only acceptable if secondary users can tell, quickly and
reliably, whether the band is free. This is the spectrum sensing problem, and it is
a hard one: the two kinds of error pull in opposite directions. A missed detection
means an IoT device transmits on top of the licensed user and causes interference; a
false alarm means free spectrum goes unused. A sensing scheme is therefore judged by
its detection probability $P_d$ at a fixed false-alarm probability $P_{fa}$.

Several sensing methods exist; energy detection is a simple non-coherent
method: it needs no knowledge of the primary signal's
waveform, timing, or carrier. The node simply measures the energy of the received
samples,
\[
T = \frac{1}{N}\sum_{n=1}^{N} |y[n]|^2,
\]
and declares the band busy if $T$ exceeds a threshold. The threshold is derived from
the noise-only distribution so that $P_{fa}$ stays at a chosen value; for large $N$
the statistic is approximately Gaussian, which yields closed-form theoretical curves
to validate against.

Energy detection has known weaknesses. Its performance collapses at low SNR, and
uncertainty in the noise power creates an ``SNR wall'' below which reliable
detection is impossible no matter how many samples are taken. A single node is also
vulnerable to deep fading and shadowing. Cooperative sensing mitigates exactly this:
several nodes sense independently and report one-bit decisions to a fusion center.
Under the OR rule, the band is declared busy if any node says busy, so a signal only
has to be caught by one of the nodes - independent nodes rarely all miss it. The
price is a higher overall false-alarm rate, which grows with the number of nodes.

In this project we build a complete simulation of this setting in Python. The
primary signal is BPSK in additive white Gaussian noise, each sensing interval uses
$N = 1000$ samples, and the threshold is set for $P_{fa} = 0.1$. We estimate
detection performance by Monte Carlo (2000 trials per point, SNR from $-20$ to
$+5$~dB), overlay the closed-form theory to validate the model, and compare a single
node with three- and five-node cooperation under the OR rule using $P_d$ vs.\ SNR
and ROC curves. The simulation code, the full numerical results, and the script
that regenerates every plot accompany this document in the project repository.

\end{document}
